\documentclass[12pt,a4paper]{article}
\usepackage[utf8]{inputenc}
\usepackage[T1]{fontenc}
\usepackage[italian]{babel}
\usepackage{amsmath,amssymb}
\usepackage{geometry}
\usepackage{hyperref}
\usepackage{enumitem}
\usepackage{booktabs}
\usepackage{array}
\usepackage{xcolor}
\usepackage{tabularx}
\usepackage{tcolorbox}

\geometry{top=2.5cm, bottom=2.5cm, left=3cm, right=3cm}
\hypersetup{colorlinks=true, linkcolor=blue!70!black}
\tcbuselibrary{skins,breakable}

\newtcolorbox{domandabox}[1]{
  colback=blue!4, colframe=blue!55!black,
  title=\textbf{Domanda: #1}, fonttitle=\bfseries\small,
  breakable, left=5pt, right=5pt, top=4pt, bottom=4pt
}
\newtcolorbox{formulabox}[1]{
  colback=yellow!8, colframe=yellow!65!black,
  title=\textbf{#1}, fonttitle=\bfseries\small,
  breakable, left=5pt, right=5pt, top=4pt, bottom=4pt
}
\newtcolorbox{attenzione}[1]{
  colback=orange!6, colframe=orange!65!black,
  title=\textbf{Attenzione: #1}, fonttitle=\bfseries\small,
  breakable, left=5pt, right=5pt, top=4pt, bottom=4pt
}
\newtcolorbox{esempiobox}{
  colback=teal!5, colframe=teal!55!black,
  title=\textbf{Esempio}, fonttitle=\bfseries\small,
  breakable, left=5pt, right=5pt, top=4pt, bottom=4pt
}
\newtcolorbox{ricordabox}{
  colback=pink!8, colframe=red!55!black,
  title=\textbf{Come ricordarlo}, fonttitle=\bfseries\small,
  breakable, left=5pt, right=5pt, top=4pt, bottom=4pt
}

\setlength{\parindent}{0pt}
\setlength{\parskip}{7pt}

\title{\textbf{Risposte Discorsive -- Sistemi Intelligenti}\\[5pt]
{\large Come rispondere all'orale e alle domande aperte}}
\author{}
\date{A.A. 2025--2026}

\begin{document}
\maketitle
\tableofcontents
\newpage

% =====================================================================
\section{Glossario degli Acronimi}
% =====================================================================

{\small
\renewcommand{\arraystretch}{1.4}
\setlength{\tabcolsep}{5pt}

\begin{tcolorbox}[colback=blue!3, colframe=blue!45!black,
  title=\textbf{Generali e Logica}, fonttitle=\bfseries\small, breakable]
\begin{tabularx}{\linewidth}{@{}p{1.5cm} p{4.6cm} X@{}}
\toprule
\textbf{Sigla} & \textbf{Significato} & \textbf{Spiegazione rapida} \\
\midrule
\textbf{IA}  & Intelligenza Artificiale    & Campo di studio: sistemi che esibiscono comportamento intelligente. \\
\textbf{ML}  & Machine Learning            & Apprendimento automatico: il sistema impara dai dati senza regole esplicite. \\
\textbf{KB}  & Knowledge Base              & Base di conoscenza: insieme di fatti e regole in forma logica. \\
\textbf{FOL} & First-Order Logic           & Logica del Primo Ordine: logica con variabili, quantificatori ($\forall$, $\exists$) e funzioni. \\
\textbf{CNF} & Conjunctive Normal Form     & Forma Normale Congiuntiva: congiunzione ($\land$) di clausole, ognuna disgiunzione ($\lor$) di letterali. Necessaria per la risoluzione. \\
\textbf{CSP} & Constraint Satisfaction Problem & Problema di Soddisfacimento di Vincoli: trovare un assegnamento completo e consistente per tutte le variabili. \\
\bottomrule
\end{tabularx}
\end{tcolorbox}

\vspace{6pt}

\begin{tcolorbox}[colback=green!3, colframe=green!50!black,
  title=\textbf{Algoritmi di Ricerca}, fonttitle=\bfseries\small, breakable]
\begin{tabularx}{\linewidth}{@{}p{1.5cm} p{4.6cm} X@{}}
\toprule
\textbf{Sigla} & \textbf{Significato} & \textbf{Spiegazione rapida} \\
\midrule
\textbf{BFS}   & Breadth-First Search        & Ricerca in ampiezza: livello per livello. Completa e ottimale (costi uguali), spazio $O(b^d)$. \\
\textbf{DFS}   & Depth-First Search          & Ricerca in profondità: esplora un ramo fino in fondo. Spazio $O(bm)$, ma non completa né ottimale. \\
\textbf{IDDFS} & Iterative Deepening DFS     & DFS con limite di profondità crescente. Unisce la memoria di DFS ($O(bd)$) e la completezza/ottimalità di BFS. \\
\textbf{UCS}   & Uniform Cost Search         & Ricerca a costo uniforme: frontiera ordinata per $g(n)$. Ottimale con costi variabili non negativi. \\
\textbf{RBFS}  & Recursive Best-First Search & A* a memoria lineare $O(bd)$: può riesaminare nodi ma non mantiene tutta la frontiera in memoria. \\
\textbf{A*}    & A-star                      & Ricerca con $f(n)=g(n)+h(n)$. Ottimale se $h$ è ammissibile (alberi) o consistente (grafi). \\
\bottomrule
\end{tabularx}
\end{tcolorbox}

\vspace{6pt}

\begin{tcolorbox}[colback=orange!3, colframe=orange!60!black,
  title=\textbf{CSP e Inferenza Logica}, fonttitle=\bfseries\small, breakable]
\begin{tabularx}{\linewidth}{@{}p{1.5cm} p{4.6cm} X@{}}
\toprule
\textbf{Sigla} & \textbf{Significato} & \textbf{Spiegazione rapida} \\
\midrule
\textbf{MRV}  & Minimum Remaining Values    & Euristica CSP ``fail-first'': sceglie la variabile col dominio residuo più piccolo. \\
\textbf{AC-3} & Arc Consistency 3           & Algoritmo per l'arc consistency: elimina valori senza supporto e propaga le riduzioni a cascata. \\
\textbf{MPG}  & Modus Ponens Generalizzato  & Modus ponens per FOL: usa l'unificazione per applicare regole universali a fatti specifici. \\
\textbf{MGU}  & Most General Unifier        & Unificatore più generale: sostituzione $\theta$ che unifica due espressioni nel modo meno restrittivo possibile. \\
\textbf{RDF}  & Resource Description Framework & Standard Web Semantico: conoscenza come triple $\langle$soggetto, predicato, oggetto$\rangle$. \\
\bottomrule
\end{tabularx}
\end{tcolorbox}

\vspace{6pt}

\begin{tcolorbox}[colback=purple!3, colframe=purple!50!black,
  title=\textbf{Machine Learning e Matrice di Confusione}, fonttitle=\bfseries\small, breakable]
\begin{tabularx}{\linewidth}{@{}p{1.5cm} p{4.6cm} X@{}}
\toprule
\textbf{Sigla} & \textbf{Significato} & \textbf{Spiegazione rapida} \\
\midrule
\textbf{TP}   & True Positive               & Vero positivo: istanza positiva classificata correttamente come positiva. \\
\textbf{TN}   & True Negative               & Vero negativo: istanza negativa classificata correttamente come negativa. \\
\textbf{FP}   & False Positive              & Falso positivo: istanza negativa classificata erroneamente come positiva. \\
\textbf{FN}   & False Negative              & Falso negativo: istanza positiva classificata erroneamente come negativa. \\
\textbf{K-NN} & K-Nearest Neighbours        & Classificatore lazy: non costruisce un modello, usa i $k$ esempi più vicini nel training set. \\
\textbf{XOR}  & eXclusive OR                & OR esclusivo: vero solo se i due input differiscono. Il perceptron non riesce ad impararlo perché non è linearmente separabile. \\
\textbf{ID3}  & Iterative Dichotomiser 3    & Algoritmo classico per alberi di decisione (variante dell'algoritmo di Hunt) basato sull'information gain. \\
\bottomrule
\end{tabularx}
\end{tcolorbox}
} % fine \small

\newpage

% =====================================================================
\section{Test di Turing e IA Forte vs Debole}
% =====================================================================

\begin{domandabox}{Cos'è il Test di Turing? Cosa dice esattamente?}
\end{domandabox}

Il Test di Turing, proposto da Alan Turing nel 1950 nel suo articolo ``Computing Machinery and Intelligence'', è un criterio operativo per stabilire se una macchina può essere considerata intelligente. L'idea è la seguente: un giudice umano interagisce via tastiera con due interlocutori che non può vedere direttamente, uno dei quali è un essere umano e l'altro è una macchina. Se il giudice non riesce a distinguere l'uno dall'altro dopo una conversazione libera, allora la macchina supera il test e può essere considerata intelligente.

La cosa importante da sottolineare è che il test non richiede che la macchina risponda \emph{correttamente} alle domande, né che le \emph{comprenda} davvero: richiede che \textbf{imiti il comportamento umano} nel rispondere. Questo è un punto critico, perché un essere umano può sbagliare, può essere distratto, può rispondere in modo incoerente, e la macchina deve essere in grado di fare altrettanto in modo convincente.

Le capacità che un sistema dovrebbe avere per superare il test riguardano l'elaborazione del linguaggio naturale, la rappresentazione della conoscenza, il ragionamento automatico e l'apprendimento. John Searle ha criticato il test con l'argomento della Stanza Cinese: una persona che segue regole per manipolare simboli cinesi senza capirne il significato starebbe ``superando'' il test senza alcuna vera comprensione. Questo ha aperto il dibattito tra IA forte e IA debole.

\begin{esempiobox}
ELIZA (1966) era un programma che simulava uno psicoterapeuta riformulando le frasi dell'utente in domande: ``Ho mal di testa'' $\to$ ``Perché dici di avere mal di testa?''. Molti utenti si convincevano di parlare con un essere umano -- un precursore del Test di Turing. Non capiva nulla: seguiva solo pattern di testo.
\end{esempiobox}

\begin{ricordabox}
\textbf{Turing = imita, non capisce.} Il test misura il comportamento esterno, non l'intelligenza interna. Una fotocopia di un libro non ``capisce'' il libro -- ma supera il test se la gente non sa distinguerla dall'originale.
\end{ricordabox}

\vspace{6pt}
\begin{domandabox}{Qual è la differenza tra IA Forte e IA Debole?}
\end{domandabox}

Quando parliamo di IA debole intendiamo sistemi che \emph{simulano} comportamenti intelligenti per uno scopo specifico, senza possedere comprensione reale o stati mentali propri. Un motore di scacchi, un sistema di riconoscimento vocale, un chatbot: questi sono tutti esempi di IA debole perché risolvono problemi definiti, ma non ``capiscono'' quello che stanno facendo nel senso in cui lo fa un essere umano.

L'IA forte, al contrario, ipotizza che una macchina possa possedere vera intelligenza, coscienza e stati mentali equivalenti a quelli umani. Non si tratta di simulare, ma di \emph{essere} intelligente. È ancora un obiettivo teorico e filosofico, non raggiunto. Il libro di testo adotta l'approccio \textbf{``agire razionalmente''}: costruire agenti che massimizzano la misura di prestazione attesa, indipendentemente dal fatto che lo facciano ``come un umano''.

\begin{ricordabox}
\textbf{Debole = fa una cosa bene (ma non sa perché).} Deep Blue batte Kasparov a scacchi $\to$ IA debole: è bravissimo a scacchi ma non sa allacciare le scarpe. \textbf{Forte = capisce, ragiona su tutto come un umano} $\to$ ancora fantascienza.
\end{ricordabox}

% =====================================================================
\section{Agenti e Ricerca}
% =====================================================================

\begin{domandabox}{Cos'è un agente razionale? Quali sono le componenti di un problema di ricerca?}
\end{domandabox}

Un agente è qualsiasi entità capace di percepire l'ambiente attraverso sensori e di agire su di esso tramite attuatori. Un agente è \emph{razionale} quando sceglie l'azione che massimizza le sue prestazioni attese sulla base della sequenza di percezioni ricevute, della conoscenza disponibile e delle azioni possibili. È importante distinguere razionalità da onniscienza: un agente razionale fa del meglio che può con le informazioni che ha, ma non è infallibile perché il mondo contiene sempre fattori ignoti.

Quando formuliamo un problema di ricerca dobbiamo definire cinque elementi fondamentali. Lo \textbf{stato iniziale} è la configurazione di partenza dell'agente nel mondo. La \textbf{funzione successore} (o modello di transizione) descrive gli effetti di ogni azione: dato uno stato $s$ e un'azione $a$, restituisce il nuovo stato risultante. Il \textbf{test obiettivo} determina se uno stato è una soluzione. Il \textbf{costo del cammino} è la somma dei costi di ogni passo compiuto. Infine, una \textbf{soluzione} è la sequenza di azioni che porta dallo stato iniziale a uno stato obiettivo.

Questi elementi non vanno confusi con un CSP: in un problema di ricerca cerchiamo una \emph{sequenza di azioni}, mentre in un CSP cerchiamo un \emph{assegnamento} che soddisfi dei vincoli, senza la nozione di cammino o costo.

\begin{esempiobox}
\textbf{8-puzzle:} stato iniziale = configurazione iniziale delle tessere; funzione successore = sposta la tessera vuota su/giù/sinistra/destra; test obiettivo = tessere in ordine 1--8; costo = numero di mosse. La soluzione è la \emph{sequenza} di spostamenti, non un singolo valore.
\end{esempiobox}

\begin{ricordabox}
\textbf{Ricerca = ``come ci arrivo?''} (sequenza di azioni). \textbf{CSP = ``cosa vale ogni variabile?''} (assegnamento). Nel problema dell'8-puzzle cerchi le mosse; nel Sudoku cerchi i numeri nelle celle.
\end{ricordabox}

% =====================================================================
\section{Algoritmi di Ricerca}
% =====================================================================

\begin{domandabox}{Descrivi BFS, DFS, IDDFS e UCS: completezza, ottimalità, complessità}
\end{domandabox}

Il \textbf{BFS} (Breadth-First Search) esplora il grafo per livelli: espande prima tutti i nodi a profondità $d$, poi quelli a profondità $d+1$, e così via. Usa una coda FIFO come frontiera. È \textbf{completo} -- trova sempre una soluzione se esiste -- e \textbf{ottimale} quando i costi degli archi sono tutti uguali, perché il primo cammino trovato è anche il più corto in numero di passi. Il difetto principale è la \textbf{complessità spaziale}: mantiene in memoria tutti i nodi del livello corrente, il che con fattore di ramificazione $b$ e soluzione a profondità $d$ richiede $O(b^d)$ nodi. Questo lo rende impraticabile per problemi con grandi spazi degli stati.

Il \textbf{DFS} (Depth-First Search) esplora in profondità lungo un ramo fino alla foglia, poi torna indietro e prova il successivo. Usa uno stack. Ha complessità spaziale $O(bm)$, dove $m$ è la profondità massima -- molto più efficiente del BFS. Tuttavia \textbf{non è completo} su grafi con cicli (può girare indefinitamente) e \textbf{non è ottimale}: può trovare soluzioni molto profonde prima di scoprirne di più vicine.

L'\textbf{IDDFS} (Iterative Deepening DFS) combina i vantaggi di BFS e DFS nel modo più elegante possibile: esegue ripetutamente una ricerca DFS limitata a profondità crescente $l = 0, 1, 2, \ldots$ fino a trovare la soluzione. Il vantaggio è che ha la \textbf{completezza e l'ottimalità di BFS} (con costi uniformi) ma la \textbf{complessità spaziale di DFS}: $O(bd)$. Il costo di riesaminare i nodi già visitati nelle iterazioni precedenti è sorprendentemente basso -- asintoticamente non cambia la complessità temporale $O(b^d)$ -- ed è di gran lunga preferibile al BFS per grandi spazi degli stati.

L'\textbf{UCS} (Uniform Cost Search) è BFS con priorità: la frontiera è una coda a priorità ordinata per $g(n)$, il costo del cammino percorso. Estrae sempre il nodo di costo minimo. È \textbf{completo e ottimale} anche con archi a costo variabile, purché i costi siano non negativi. Quando tutti i costi sono uguali a 1, degenera esattamente in BFS.

\begin{esempiobox}
Immagina un labirinto con corridoi di lunghezza diversa. \textbf{BFS} trova l'uscita con meno svolte; \textbf{DFS} imbocca il primo corridoio e va fino in fondo (può sbagliare tutto); \textbf{IDDFS} prova prima entro 1 corridoio, poi 2, poi 3\ldots\ trovando la soluzione più breve; \textbf{UCS} tiene conto della lunghezza reale di ogni corridoio e trova il percorso più corto in metri.
\end{esempiobox}

\begin{ricordabox}
\textbf{Memoria}: BFS = \textbf{B}ello ma \textbf{B}ulimica di RAM. DFS = \textbf{D}iscende \textbf{D}ritto ma si perde. IDDFS = il migliore dei due mondi. UCS = BFS con bilancia (pesa i costi).\\[3pt]
\textbf{Tabella da ricordare}: DFS è l'unico \emph{non} completo e \emph{non} ottimale. IDDFS è completo+ottimale con spazio di DFS. UCS è completo+ottimale con costi variabili.
\end{ricordabox}

\begin{center}
\renewcommand{\arraystretch}{1.3}
\begin{tabular}{lcccc}
\toprule
\textbf{Algoritmo} & \textbf{Completo} & \textbf{Ottimale} & \textbf{Tempo} & \textbf{Spazio} \\
\midrule
BFS       & Sì & Sì (costi uguali) & $O(b^d)$   & $O(b^d)$ \\
DFS       & No & No                & $O(b^m)$   & $O(bm)$ \\
IDDFS     & Sì & Sì (costi uguali) & $O(b^d)$   & $O(bd)$ \\
UCS       & Sì & Sì                & $O(b^{C^*/\varepsilon})$ & $O(b^{C^*/\varepsilon})$ \\
A*        & Sì & Sì (euristica ammiss.) & $O(b^d)$ & $O(b^d)$ \\
\bottomrule
\end{tabular}
\end{center}

\begin{domandabox}{Cosa differenzia un algoritmo informato da uno non informato?}
\end{domandabox}

La differenza fondamentale è l'utilizzo di una \textbf{funzione euristica} $h(n)$. Un algoritmo non informato come BFS, DFS o UCS non sa nulla sulla posizione del goal rispetto allo stato corrente: esplora il grafo in modo cieco, seguendo solo la struttura dello spazio degli stati. Un algoritmo informato, come A* o la ricerca greedy, dispone invece di una stima $h(n)$ del costo per raggiungere il goal da $n$ e la usa per guidare l'esplorazione verso le zone più promettenti.

È importante chiarire che la funzione costo $g(n)$ -- il costo del cammino percorso finora -- è usata anche dagli algoritmi non informati come UCS. Quello che distingue gli algoritmi informati non è quindi il costo, ma proprio la stima euristica del costo futuro.

\begin{domandabox}{Cos'è un'euristica ammissibile? E una dominante?}
\end{domandabox}

Un'euristica $h(n)$ è \textbf{ammissibile} se non sovrastima mai il costo reale per raggiungere il goal: formalmente $h(n) \leq h^*(n)$ per ogni nodo $n$, dove $h^*(n)$ è il costo ottimo reale. Questo la rende ``ottimistica'': preferisce sempre soluzioni buone. La proprietà fondamentale è che A* con un'euristica ammissibile garantisce l'ottimalità su alberi di ricerca.

Se però lavoriamo su grafi con una lista chiusa -- dove i nodi già espansi non vengono riesaminati -- abbiamo bisogno di una condizione più forte: la \textbf{consistenza} (o monotonicità). Un'euristica è consistente se $h(n) \leq c(n, n') + h(n')$ per ogni arco $(n, n')$. La consistenza implica l'ammissibilità, ma non viceversa. Un'euristica consistente garantisce l'ottimalità di A* anche su grafi.

Un'euristica $h_1$ \textbf{domina} $h_2$ se $h_1(n) \geq h_2(n)$ per \emph{ogni} nodo $n$. Questo significa che $h_1$ è più informativa: fa stime più alte (senza mai superare il valore reale), e quindi A* con $h_1$ espande meno nodi rispetto ad A* con $h_2$.

Una funzione costante zero, $h(n) = 0$, è tecnicamente ammissibile ma completamente non informativa: A* con questa euristica degenera in UCS e non guida in alcun modo la ricerca verso il goal.

\begin{esempiobox}
Nell'\textbf{8-puzzle}: $h_1$ = numero di tessere fuori posto (ammissibile: non sovrastima mai le mosse necessarie). $h_2$ = distanza di Manhattan di ogni tessera dalla sua posizione finale (ammissibile e \emph{dominante} su $h_1$: stima più alta ma ancora ottimistica). A* con $h_2$ espande molti meno nodi che con $h_1$.
\end{esempiobox}

\begin{ricordabox}
\textbf{Ammissibile = ottimista = non mente mai per eccesso.} Come un GPS che dice ``arrivi in 10 minuti'': può sbagliare per difetto (arrivare in 12) ma non per eccesso.\\[3pt]
\textbf{Consistente = monotona}: $f(n)$ non scende mai lungo il cammino. Se la stima cambia ogni volta che aggiungi un passo, sei inconsistente.\\[3pt]
\textbf{Domina}: $h_1 \geq h_2$ ovunque $\Rightarrow$ $h_1$ sa di più $\Rightarrow$ A* con $h_1$ è più furbo.
\end{ricordabox}

\begin{domandabox}{Come funziona A*? Perché è ottimale?}
\end{domandabox}

A* è un algoritmo di ricerca informata che ordina la frontiera in base alla funzione $f(n) = g(n) + h(n)$, dove $g(n)$ è il costo del cammino già percorso da radice a $n$ e $h(n)$ è la stima euristica del costo rimanente per raggiungere il goal. A ogni passo espande il nodo con $f$ minimo.

L'ottimalità di A* dipende dalla proprietà dell'euristica. Su \textbf{alberi di ricerca}, basta l'ammissibilità: $h(n) \leq h^*(n)$. Su \textbf{grafi con lista chiusa}, dove i nodi espansi non vengono rivisitati, serve la \textbf{consistenza}: $h(n) \leq c(n,n') + h(n')$ per ogni arco. La consistenza garantisce che $f$ non diminuisca mai lungo un cammino, ovvero che il primo percorso con cui si espande un nodo sia già quello ottimale.

La differenza tra A* e UCS è proprio questa: UCS usa solo $g(n)$ e non sfrutta nessuna conoscenza del goal, mentre A* usa anche $h(n)$ per orientare la ricerca. Con $h = 0$, A* degenera in UCS.

\begin{domandabox}{Cos'è RBFS e quando si usa?}
\end{domandabox}

RBFS (Recursive Best-First Search) è una versione di A* progettata per operare con \textbf{memoria lineare}: mantiene in memoria solo il cammino corrente dalla radice al nodo attivo e i valori $f$ dei nodi fratelli non ancora esplorati. Quando un ramo supera una soglia $f_{limit}$, l'algoritmo torna indietro, aggiorna il valore del nodo con il miglior $f$ trovato nel sottoalbero esplorato, e poi ricomincia dal ramo con $f$ minore.

Rispetto ad A*, RBFS può riesaminare nodi più volte a causa dei ritorni all'indietro, aumentando il costo temporale; però la complessità spaziale è solo $O(bd)$ -- analoga a IDDFS. È l'algoritmo da preferire quando la memoria disponibile è il fattore limitante.

\begin{domandabox}{Cos'è la ricerca bidirezionale?}
\end{domandabox}

La ricerca bidirezionale esegue simultaneamente \textbf{due ricerche in parallelo}: una in avanti dallo stato iniziale e una all'indietro dallo stato obiettivo, fermandosi quando le due frontiere si incontrano. L'idea è che invece di esplorare un albero di profondità $d$ con complessità $O(b^d)$, si esplorano due alberi di profondità $d/2$ ciascuno, riducendo la complessità a $O(b^{d/2})$ -- un risparmio esponenziale.

La difficoltà pratica è definire come effettuare la ricerca all'indietro: bisogna avere azioni invertibili e sapere in anticipo lo stato obiettivo (o essere in grado di generare i predecessori di uno stato).

% =====================================================================
\section{CSP -- Problemi di Soddisfacimento di Vincoli}
% =====================================================================

\begin{domandabox}{Cos'è un CSP? Come si definisce una soluzione?}
\end{domandabox}

Un problema di soddisfacimento di vincoli (CSP) è una tripla $\langle X, D, C \rangle$ dove $X$ è un insieme di variabili, $D$ è l'insieme dei rispettivi domini e $C$ è l'insieme dei vincoli. Ogni vincolo specifica quali combinazioni di valori sono ammissibili per un sottoinsieme di variabili.

Una soluzione di un CSP è un \textbf{assegnamento completo e consistente}: completo perché ogni variabile deve avere un valore, consistente perché nessun vincolo deve essere violato. Questo lo distingue nettamente dalla ricerca classica, dove la soluzione è una sequenza di azioni. Nel CSP non abbiamo né un cammino né un costo: vogliamo solo trovare una configurazione valida.

\begin{esempiobox}
\textbf{Colorazione mappa Australia}: variabili $X = \{WA, NT, SA, Q, NSW, V, T\}$, domini $D = \{Rosso, Verde, Blu\}$, vincoli: ogni coppia di regioni adiacenti deve avere colori diversi (es. $WA \neq NT$, $WA \neq SA$, \ldots). Una soluzione: $WA=R, NT=G, SA=B, Q=R, NSW=G, V=R, T=G$.
\end{esempiobox}

\begin{ricordabox}
\textbf{CSP = Sudoku}: hai celle (variabili), valori 1--9 (domini), regole di riga/colonna/quadrato (vincoli). Non cerchi come arrivare alla soluzione -- cerchi i valori giusti.
\end{ricordabox}

\begin{domandabox}{Cos'è la commutatività nel CSP? Perché riduce l'ampiezza dell'albero?}
\end{domandabox}

La commutatività nei CSP afferma che l'\textbf{ordine in cui si assegnano le variabili non influenza la soluzione}: se $\{X_1 = a, X_2 = b\}$ è una soluzione, lo è indipendentemente dal fatto che si sia assegnato prima $X_1$ o prima $X_2$.

Questa proprietà è cruciale perché permette di costruire l'albero di ricerca assegnando \emph{una variabile alla volta} invece di considerare tutte le possibili sequenze di assegnamento. Senza sfruttare la commutatività, l'albero avrebbe $n!$ percorsi distinti per ogni soluzione (uno per ogni permutazione dell'ordine); con la commutatività se ne considera uno solo. Questo riduce drasticamente l'\textbf{ampiezza} dell'albero (il fattore di ramificazione passa da $n \cdot d$ a $d$).

\begin{ricordabox}
\textbf{Sudoku}: assegnare prima la cella A1 e poi B1 è lo stesso che assegnare B1 e poi A1. L'ordine in cui compili le celle non cambia la soluzione finale -- ma scegliere \emph{quale} cella compilare prima (MRV) può risparmiarti un sacco di tentativi.
\end{ricordabox}

\begin{domandabox}{Quali algoritmi si usano per risolvere un CSP? Qual è il più e il meno efficiente?}
\end{domandabox}

L'approccio più basilare -- e meno efficiente -- è il \textbf{generate and test}: si generano tutti i possibili assegnamenti completi dell'intero set di variabili, e per ciascuno si verifica se i vincoli sono soddisfatti. La complessità è $O(d^n)$ dove $d$ è la dimensione del dominio e $n$ è il numero di variabili, perché non sfrutta i vincoli durante la costruzione.

Il \textbf{backtracking} migliora notevolmente questo approccio: si assegnano le variabili una alla volta e, non appena un assegnamento parziale viola un vincolo, si torna indietro senza esplorare tutte le sue estensioni. Sfruttare la commutatività degli assegnamenti riduce ulteriormente l'ampiezza dell'albero.

L'algoritmo più efficiente tra quelli tipicamente discussi è il \textbf{back-jumping}. Mentre il backtracking cronologico, in caso di fallimento, torna sempre alla variabile assegnata immediatamente prima (anche se non è la responsabile del conflitto), il back-jumping identifica il \textbf{conflict set} -- l'insieme delle variabili già assegnate che hanno causato il fallimento -- e torna direttamente alla più recente tra queste. Questo evita di esplorare rami che non avrebbero comunque portato a una soluzione.

\begin{domandabox}{Cosa sono MRV e l'euristica di grado nel CSP?}
\end{domandabox}

Queste sono due euristiche per scegliere in quale ordine assegnare le variabili durante il backtracking, con l'obiettivo di rilevare i fallimenti il prima possibile.

La \textbf{MRV} (Minimum Remaining Values), detta anche euristica ``fail-first'', sceglie la variabile con il \textbf{dominio residuo più piccolo}: si preferisce la variabile che ha meno valori ammissibili rimanenti, perché è quella che con più probabilità porterà a un fallimento. Rilevare il fallimento presto significa non sprecare tempo a espandere rami inutili.

L'\textbf{euristica di grado} (degree heuristic) serve come criterio di spareggio quando due o più variabili hanno lo stesso numero di valori residui: si sceglie la variabile coinvolta nel maggior numero di vincoli con le variabili ancora non assegnate. L'idea è che la variabile con più connessioni avrà un impatto maggiore sulle scelte future, quindi è meglio vincolarla subito.

\begin{esempiobox}
Australia: $SA$ (South Australia) tocca 5 altre regioni. Se assegni prima $SA$, riduci immediatamente le scelte di tutte le sue vicine. MRV direbbe: scegli chi ha meno colori rimasti; euristica di grado dice: a parità, scegli $SA$ perché ha più vincoli.
\end{esempiobox}

\begin{ricordabox}
\textbf{MRV = fail-first}: scegli chi rischia di bloccarti prima, così fallisci subito e risparmi lavoro inutile. \textbf{Euristica di grado = chi comanda}: scegli chi ha più vincoli con gli altri = chi influenza di più il resto.
\end{ricordabox}

\begin{domandabox}{Cos'è il Forward Checking nel CSP?}
\end{domandabox}

Il forward checking è una tecnica di \textbf{propagazione dei vincoli} che opera \emph{durante} la costruzione dell'assegnamento: ogni volta che si assegna un valore alla variabile $X_i$, si controllano immediatamente tutti i vincoli che coinvolgono $X_i$ e una variabile ancora non assegnata $X_j$, eliminando dal dominio di $X_j$ i valori incompatibili con il valore appena scelto per $X_i$.

Il vantaggio è che questo rileva fallimenti molto prima del backtracking puro: se il dominio di una variabile futura diventa vuoto, si può già fare backtracking senza dover assegnare altre variabili. Tuttavia il forward checking non è completo nel propagare le inconsistenze: non rileva le incompatibilità tra variabili future non ancora assegnate. Per questo, una tecnica più potente è l'\textbf{arc consistency}.

\begin{domandabox}{Cos'è l'Arc Consistency (AC-3)?}
\end{domandabox}

Un arco $(X_i, X_j)$ in un CSP è \textbf{arc-consistent} se per ogni valore nel dominio di $X_i$ esiste almeno un valore nel dominio di $X_j$ che è compatibile con il vincolo tra $X_i$ e $X_j$. L'arc consistency si ottiene eliminando iterativamente i valori che violano questa condizione.

L'algoritmo \textbf{AC-3} gestisce una coda di archi da verificare. Prende un arco $(X_i, X_j)$, rimuove da $D_i$ i valori per cui non esiste alcun supporto in $D_j$. Se $D_i$ si riduce, aggiunge alla coda tutti gli archi $(X_k, X_i)$ con $k \neq j$, perché la riduzione del dominio di $X_i$ potrebbe rendere inconsistenti archi già controllati. L'algoritmo termina quando la coda è vuota o quando un dominio diventa vuoto (inconsistenza rilevata).

La differenza con il forward checking è che AC-3 propaga le riduzioni a cascata tra tutte le coppie di variabili, non solo tra la variabile appena assegnata e le future: è quindi molto più potente nel rilevare fallimenti anticipatamente, a costo di più computazione. L'arc consistency si può applicare sia come pre-processo prima del backtracking, sia interleaved durante la ricerca.

\begin{domandabox}{Cos'è la Path Consistency?}
\end{domandabox}

La path consistency (o 3-consistenza) estende il concetto di arc consistency a \textbf{triple di variabili}. Un insieme di tre variabili $\{X_i, X_j, X_k\}$ è path-consistent se per ogni assegnamento consistente di $X_i$ e $X_j$ esiste un valore di $X_k$ che è compatibile con entrambi.

La gerarchia è: node consistency (1 variabile) $\subset$ arc consistency (2 variabili) $\subset$ path consistency (3 variabili) $\subset \ldots$ Più si sale nella gerarchia, più si rilevano inconsistenze in anticipo ma a costo computazionale crescente. Nella pratica, arc consistency è spesso il miglior compromesso.

\begin{domandabox}{Cos'è il back-jumping e il conflict set?}
\end{domandabox}

Quando una variabile $X_k$ non può ricevere alcun valore senza violare un vincolo, il suo conflict set è l'insieme delle variabili già assegnate che sono in conflitto con $X_k$. Formalmente: $CS(X_k) = \{X_j : j < k, \exists \text{ vincolo tra } X_j \text{ e } X_k\}$.

Il back-jumping salta direttamente all'ultima variabile nel conflict set, bypassando tutte le variabili intermedie che non hanno contribuito al fallimento. Questo lo rende molto più efficiente del semplice backtracking cronologico, specialmente in problemi con molte variabili e vincoli distribuiti in modo sparso.

\begin{esempiobox}
Variabili $X_1, X_2, X_3, X_4, X_5$. Assegni $X_1, X_2, X_3, X_4$ e poi $X_5$ fallisce perché in conflitto solo con $X_1$ e $X_3$. Il backtracking cronologico torna a $X_4$ (inutile). Il back-jumping salta direttamente a $X_3$ (l'ultimo nel conflict set $\{X_1, X_3\}$), risparmiando tutta l'esplorazione di $X_4$.
\end{esempiobox}

\begin{ricordabox}
\textbf{Backtracking = torna dal vicino. Back-jumping = torna dal colpevole.} Se litighi con qualcuno ma l'intermediario non c'entra, vai direttamente dal responsabile.
\end{ricordabox}

% =====================================================================
\section{Ricerca Avversariale}
% =====================================================================

\begin{domandabox}{Come funziona l'algoritmo Minimax?}
\end{domandabox}

Minimax è un algoritmo per la ricerca in problemi con avversario, tipicamente giochi a somma zero a due giocatori con informazione perfetta come scacchi, dama o tris. Un gioco a \textbf{somma zero} è tale che il guadagno di un giocatore corrisponde esattamente alla perdita dell'altro: quello che guadagna MAX viene perso da MIN.

Il funzionamento di Minimax si basa sull'assunzione che entrambi i giocatori giochino in modo ottimale. L'algoritmo costruisce l'albero di gioco espandendo le mosse possibili a partire dallo stato corrente, alternativamente per MAX e per MIN. Una volta raggiunte le foglie -- stati terminali o stati alla profondità limite -- si applica la \textbf{funzione di utilità} (per stati terminali) o la \textbf{funzione di valutazione} EVAL (per stati non terminali alla profondità limite).

Il punto cruciale è che i valori vengono poi propagati \textbf{dalle foglie verso la radice}: i nodi MAX scelgono il valore massimo tra i figli, i nodi MIN scelgono il minimo. In questo modo la radice dell'albero acquisisce il valore della mossa ottimale che MAX dovrebbe eseguire.

È importante chiarire che Minimax non costruisce un percorso che porta MAX a vincere: \textbf{suggerisce soltanto la prossima mossa ottimale}, assumendo che l'avversario risponda nel modo migliore possibile. La vittoria non è garantita se la posizione di partenza è già sfavorevole.

\begin{esempiobox}
Albero a 2 livelli: MAX è alla radice, MIN ai figli, foglie = valori di utilità.
\begin{center}
\renewcommand{\arraystretch}{1.1}
\begin{tabular}{ccc}
& \textbf{MAX} & \\
\textbf{MIN}: $\min(3,5)=3$ & & \textbf{MIN}: $\min(2,9)=2$ \\
foglie: 3, 5 & & foglie: 2, 9
\end{tabular}
\end{center}
MAX sceglie $\max(3,2)=3$: va a sinistra. MIN avrebbe risposto con il valore 3, non 5.
\end{esempiobox}

\begin{ricordabox}
\textbf{Le foglie risalgono: MIN taglia i rami alti, MAX taglia i rami bassi.} Immagina MAX come te che vuoi il voto più alto possibile, e MIN come il professore che cerca di darti il minimo. Alla fine vince il meglio raggiungibile in uno scenario in cui il prof fa del suo peggio.
\end{ricordabox}

\begin{formulabox}{Formula Minimax}
\[
\text{MiniMax}(s) =
\begin{cases}
\text{Utility}(s) & \text{se } s \text{ è terminale} \\[4pt]
\displaystyle\max_{a} \;\text{MiniMax}(\text{Result}(s,a)) & \text{se Player}(s) = \text{MAX} \\[8pt]
\displaystyle\min_{a} \;\text{MiniMax}(\text{Result}(s,a)) & \text{se Player}(s) = \text{MIN}
\end{cases}
\]
Complessità temporale: $O(b^m)$. Con alpha-beta pruning nel caso migliore: $O(b^{m/2})$.
\end{formulabox}

\begin{domandabox}{Cos'è l'Alpha-Beta Pruning? Come migliora Minimax?}
\end{domandabox}

L'alpha-beta pruning è un'ottimizzazione di Minimax che \textbf{elimina rami dell'albero di gioco che non possono influenzare la decisione finale}, senza cambiare il risultato. L'idea è mantenere due valori durante la visita: $\alpha$ (il miglior valore garantito finora per MAX) e $\beta$ (il miglior valore garantito finora per MIN, ovvero il più basso).

La regola di potatura è la seguente: se durante l'esplorazione di un nodo MIN troviamo un valore $\leq \alpha$, possiamo smettere di espandere quel nodo (potatura beta), perché MAX non sceglierà mai questo ramo -- ha già un'alternativa migliore. Simmetricamente, se durante un nodo MAX troviamo un valore $\geq \beta$, possiamo tagliare (potatura alfa), perché MIN non permetterà mai a MAX di arrivare qui.

Nel \textbf{caso migliore} -- quando i nodi sono ordinati in modo ottimale, con le mosse migliori esaminate per prime -- l'alpha-beta dimezza la profondità effettiva dell'albero: la complessità scende da $O(b^m)$ a $O(b^{m/2})$, il che equivale a poter cercare al doppio della profondità con la stessa quantità di calcolo. Nel \textbf{caso peggiore} (ordinamento pessimo) non dà nessun vantaggio.

È fondamentale capire che alpha-beta \emph{non cambia la mossa scelta da Minimax}: produce esattamente lo stesso risultato, solo esplorando meno nodi. La potatura non influenza la correttezza perché i rami tagliati non avrebbero mai potuto essere scelti da un giocatore razionale.

\begin{ricordabox}
\textbf{$\alpha$ = il meglio garantito a MAX finora} (cresce). \textbf{$\beta$ = il meglio garantito a MIN finora} (decresce). Taglia quando $\alpha \geq \beta$: MAX non sceglierà mai un ramo dove MIN può già fare di meglio altrove.\\[3pt]
\textbf{Analogia}: stai comprando un'auto (MAX vuoi spendere poco, MIN = concessionario vuole guadagnare). Hai già un'offerta a 10k€. Se trovi un concessionario che parte da 15k€ senza sconti, non perdi tempo a trattare: tagli quel ramo.
\end{ricordabox}

% =====================================================================
\section{Logica del Primo Ordine}
% =====================================================================

\begin{domandabox}{Quali simboli della FOL richiedono interpretazione?}
\end{domandabox}

In FOL distinguiamo simboli logici (quantificatori, connettivi, uguaglianza) che hanno un significato fisso, e simboli \textbf{non logici} che richiedono un'interpretazione, ovvero una mappatura verso il dominio del discorso.

I simboli non logici che richiedono interpretazione sono tre: i \textbf{simboli di costante} (denotano un oggetto specifico del dominio), i \textbf{simboli di funzione} (denotano una funzione dal dominio a sé stesso) e i \textbf{simboli di predicato} (denotano una relazione sul dominio). Le variabili non richiedono interpretazione perché sono legate dai quantificatori: il loro significato dipende dall'assegnamento corrente, non dall'interpretazione del modello.

\begin{domandabox}{Cos'è la Forma Normale Congiuntiva (CNF)?}
\end{domandabox}

La Forma Normale Congiuntiva è una rappresentazione standard delle formule logiche come \textbf{congiunzione di clausole}, dove ogni clausola è una \textbf{disgiunzione di letterali}. Un letterale è un atomo o la sua negazione. Per esempio: $(A \lor \neg B) \land (B \lor C \lor \neg A)$ è in CNF.

La CNF è fondamentale perché è la forma richiesta dall'algoritmo di \textbf{risoluzione}. Qualsiasi formula della logica proposizionale o del primo ordine può essere convertita in CNF attraverso una procedura standard: si eliminano prima i bicondizionali ($\Leftrightarrow$) e le implicazioni ($\Rightarrow$), poi si porta la negazione verso l'interno usando le leggi di De Morgan, e infine si distribuisce la disgiunzione sulla congiunzione.

\begin{domandabox}{Cosa sono le clausole di Horn e perché sono importanti?}
\end{domandabox}

Una clausola di Horn è una clausola che contiene \textbf{al più un letterale positivo} (non negato). Per esempio, $\neg P(x) \lor Q(x) \lor \neg R(x)$ è di Horn perché ha un solo letterale positivo ($Q(x)$), mentre $P(x) \lor Q(x) \lor \neg R(x)$ \emph{non} è di Horn perché ha due letterali positivi.

Le clausole di Horn sono importanti per due ragioni principali. Prima di tutto, l'inferenza su insiemi di clausole di Horn è decidibile in \textbf{tempo polinomiale}, mentre l'inferenza sulla logica del primo ordine in generale non è decidibile. Secondariamente, le clausole di Horn ammettono algoritmi di inferenza molto efficienti come il \textbf{forward chaining} e il \textbf{backward chaining}.

In forma implicativa, una clausola di Horn con un letterale positivo si scrive come $P_1 \land P_2 \land \cdots \land P_k \Rightarrow Q$: le parti negative sono le premesse, la parte positiva è la conclusione. Questo corrisponde esattamente alla struttura delle regole nei sistemi esperti e in Prolog.

\begin{esempiobox}
$\neg \text{Umano}(x) \lor \text{Mortale}(x)$ è di Horn (un solo positivo: Mortale). Equivale a $\text{Umano}(x) \Rightarrow \text{Mortale}(x)$.\\
$\text{Umano}(x) \lor \text{Robot}(x) \lor \neg\text{Costruito}(x)$ \emph{non} è di Horn: due letterali positivi (Umano e Robot).
\end{esempiobox}

\begin{ricordabox}
\textbf{Horn = al più 1 testa = al più 1 conclusione.} Le clausole di Horn sono le regole ``SE...ALLORA...''. Se ci fossero due teste, la regola direbbe ``allora A oppure B'' -- ambigua, non più Horn. Prolog lavora solo con clausole di Horn: ogni clausola ha esattamente una testa (o nessuna, per i fatti).
\end{ricordabox}

\begin{domandabox}{Come funziona il Forward Chaining?}
\end{domandabox}

Il forward chaining è un algoritmo di inferenza \textbf{guidato dai dati}: parte dai fatti noti nella base di conoscenza e applica ripetutamente le regole le cui premesse sono soddisfatte, aggiungendo nuove conclusioni ai fatti noti. Continua fino al raggiungimento del goal o fino al punto fisso, ovvero fino a quando non si possono più derivare nuovi fatti.

Perché il forward chaining sia applicabile, è necessario che la base di conoscenza contenga \emph{clausole di Horn} e che \textbf{tutti} i fatti nelle premesse di una regola siano già presenti in KB. Se anche uno solo dei fatti richiesti manca, la regola non può essere applicata.

Il \textbf{Modus Ponens Generalizzato} (MPG) è la versione del forward chaining per la logica del primo ordine: invece di lavorare su formule proposizionali, usa l'\textbf{unificazione} per trovare la sostituzione $\theta$ che rende le premesse della regola identiche ai fatti noti, e poi applica quella sostituzione alla conclusione.

\begin{formulabox}{Modus Ponens -- formula}
\[
\frac{\alpha \qquad \alpha \Rightarrow \beta}{\beta}
\qquad\qquad
\text{MPG: }\frac{p_1,\ldots,p_k \qquad (P_1 \land \cdots \land P_k \Rightarrow Q)\theta}{Q\theta}
\]
dove $\theta$ è la sostituzione che unifica $P_i$ con $p_i$ per ogni $i$.
\end{formulabox}

\begin{domandabox}{Cos'è il Backward Chaining? Differenze con il Forward Chaining?}
\end{domandabox}

Il backward chaining è un algoritmo di inferenza \textbf{guidato dall'obiettivo}: parte dalla query che si vuole dimostrare e cerca \emph{a ritroso} le regole la cui conclusione corrisponde alla query, tentando poi di dimostrare le premesse di quelle regole ricorsivamente. È la strategia alla base di \textbf{Prolog}.

La differenza fondamentale con il forward chaining è la direzione dell'inferenza. Il forward chaining parte dai fatti e deriva tutte le conseguenze possibili finché non si raggiunge il goal -- può derivare molti fatti inutili lungo il percorso. Il backward chaining parte dal goal e determina solo quali fatti servono per dimostrarlo -- è molto più efficiente quando la query è specifica e lo spazio degli obiettivi parziali è piccolo rispetto all'insieme di tutti i fatti derivabili.

Entrambi richiedono clausole di Horn per funzionare correttamente. La scelta tra i due dipende dal problema: il forward chaining è preferibile quando si hanno pochi fatti e molte query, mentre il backward chaining è preferibile con molti fatti e query specifiche.

\begin{esempiobox}
KB: Umano(Socrate), $\text{Umano}(x) \Rightarrow \text{Mortale}(x)$.\\
\textbf{Forward chaining}: parto da Umano(Socrate), applico la regola $\Rightarrow$ derivo Mortale(Socrate). Continuo finché non raggiungo il goal.\\
\textbf{Backward chaining}: voglio dimostrare Mortale(Socrate). Cerco regole con testa Mortale(x): trovo $\text{Umano}(x) \Rightarrow \text{Mortale}(x)$. Sostituisco $\{x/\text{Socrate}\}$ e mi chiedo: Umano(Socrate)? Sì, è in KB. Fine.
\end{esempiobox}

\begin{ricordabox}
\textbf{Forward = detective che parte dagli indizi e vede dove porta.} \textbf{Backward = avvocato che parte dalla tesi e cerca le prove.} Se la domanda è specifica (``Socrate è mortale?'') il backward è più diretto; se hai pochi fatti e vuoi tutto il deducibile, il forward è meglio.
\end{ricordabox}

\begin{domandabox}{Cos'è l'unificazione? Cosa si intende per MGU?}
\end{domandabox}

L'unificazione è il processo di trovare una \textbf{sostituzione} $\theta$ -- un'associazione di variabili a termini -- tale che due espressioni diventino sintatticamente identiche dopo aver applicato $\theta$. Per esempio, $P(x, f(y))$ e $P(a, f(b))$ si unificano con la sostituzione $\{x/a, y/b\}$: applicando $\theta$, entrambe diventano $P(a, f(b))$.

Il \textbf{Most General Unifier} (MGU) è la sostituzione più generale tra tutte quelle che unificano due espressioni: non vincola le variabili più del necessario. Se esistono più sostituzioni unificanti, l'MGU è quella da cui si possono derivare tutte le altre. Per esempio, $\{x/y\}$ è più generale di $\{x/a, y/a\}$ perché la seconda si ottiene dalla prima con un'ulteriore sostituzione.

L'unificazione fallisce quando si incontra l'\textbf{occur check}: non si può unificare $x$ con $f(x)$ perché si creerebbe una struttura infinita. In molte implementazioni pratiche (come Prolog) l'occur check viene omesso per motivi di efficienza, a rischio di risultati errati in casi particolari.

L'unificazione è il meccanismo fondamentale che rende possibile il Modus Ponens Generalizzato e la risoluzione in FOL: senza di essa non potremmo applicare regole universali a fatti specifici.

\begin{esempiobox}
UNIFY($\text{Conosce}(\text{Anna}, x)$,\; $\text{Conosce}(\text{Anna}, \text{Giovanni})$) $= \{x/\text{Giovanni}\}$ \\
UNIFY($P(x, f(y))$,\; $P(a, f(b))$) $= \{x/a,\; y/b\}$ \\
UNIFY($P(x)$,\; $P(f(x))$) $=$ \textbf{fallisce} (occur check: $x$ apparirebbe dentro $f(x)$, struttura infinita)
\end{esempiobox}

\begin{ricordabox}
\textbf{Unificazione = trovare i valori delle variabili che rendono due espressioni identiche.} È come risolvere un'equazione: $P(x, 5) = P(3, y)$ $\Rightarrow$ $x=3, y=5$. MGU = la soluzione più generale: non dire $x=3$ se basta dire $x=y$.
\end{ricordabox}

\begin{domandabox}{Cos'è la risoluzione e come si usa per la refutazione?}
\end{domandabox}

La risoluzione è una regola di inferenza che, date due clausole contenenti un letterale e il suo complementare, produce una nuova clausola -- il \emph{risolvente} -- eliminando quei letterali. Per esempio, da $(A \lor C)$ e $(\neg C \lor B)$ si ottiene il risolvente $(A \lor B)$: i letterali complementari $C$ e $\neg C$ vengono eliminati.

In FOL, si usa prima l'\textbf{unificazione} per rendere complementari i letterali prima di risolvere. A differenza del forward chaining, la risoluzione non richiede clausole di Horn: funziona su qualsiasi clausola.

La \textbf{dimostrazione per refutazione} sfrutta questo principio: per dimostrare che $KB \models \alpha$, si aggiunge $\neg\alpha$ alla base di conoscenza convertita in CNF, e si applicano ripetutamente le risoluzioni. Se si arriva alla \textbf{clausola vuota} $\square$ -- ovvero a una contraddizione -- significa che $KB \land \neg\alpha$ è insoddisfacibile, e quindi $\alpha$ è conseguenza logica di KB.

\begin{esempiobox}
KB: (1) $A \lor B$\quad (2) $\neg B \lor C$\quad Voglio dimostrare $C$ assumendo $A$ falsa.\\
Aggiungo $\neg A$ e $\neg C$ (la negazione della tesi).\\
Risolvo (1) con $\neg A$: ottengo $B$. Risolvo $B$ con (2): ottengo $C$. Risolvo $C$ con $\neg C$: ottengo $\square$. Contraddizione! Quindi $C$ è dimostrato.
\end{esempiobox}

\begin{ricordabox}
\textbf{Refutazione = dimostrazione per assurdo.} ``Supponiamo che $\alpha$ sia falso. Questo porta a contraddizione. Quindi $\alpha$ è vero.'' Clausola vuota $\square$ = hai trovato la contraddizione.
\end{ricordabox}

\begin{domandabox}{Cosa sono i termini ground e le funzioni di Skolem?}
\end{domandabox}

Un termine è \textbf{ground} se non contiene variabili: contiene solo costanti o applicazioni di funzioni a costanti. Per esempio, $\text{Padre}(\text{Giovanni})$ e $\text{Anna}$ sono termini ground, mentre $\text{Padre}(x)$ non lo è perché contiene la variabile $x$. Una formula è ground se tutti i suoi termini lo sono.

Le \textbf{funzioni di Skolem} emergono durante la conversione di formule FOL in CNF, precisamente nella fase di \emph{skolemizzazione}: l'eliminazione dei quantificatori esistenziali. Quando incontriamo $\exists x\; P(x)$ in una formula senza quantificatori universali esterni, sostituiamo $x$ con una nuova costante di Skolem $k$. Quando invece il quantificatore esistenziale è dentro uno universale -- per esempio $\forall y\; \exists x\; P(x, y)$ -- sostituiamo $x$ con una \textbf{funzione di Skolem} $f(y)$, perché il valore di $x$ può dipendere da $y$. La skolemizzazione preserva la soddisfacibilità della formula, non la validità.

\begin{esempiobox}
$\exists x\; \text{Gatto}(x)$ $\xrightarrow{\text{Skolem}}$ $\text{Gatto}(k)$ \quad (k = nome di un gatto specifico, ad es. ``Whiskers'')\\
$\forall y\; \exists x\; \text{Padre}(x, y)$ $\xrightarrow{\text{Skolem}}$ $\forall y\; \text{Padre}(f(y), y)$ \quad (f = ``il padre di'')\\
Un termine ground non ha variabili: $\text{Padre}(\text{Mario}, \text{Luigi})$ è ground; $\text{Padre}(x, \text{Luigi})$ non lo è.
\end{esempiobox}

\begin{ricordabox}
\textbf{Costante di Skolem}: ``qualcuno esiste'' $\to$ dagli un nome. \textbf{Funzione di Skolem}: ``per ogni Y esiste qualcuno'' $\to$ quel qualcuno dipende da Y $\to$ usa una funzione $f(Y)$.\\
\textbf{Ground = senza variabili} = tutto istanziato = puoi valutarlo direttamente.
\end{ricordabox}

\begin{domandabox}{Cosa sono tautologie e contraddizioni?}
\end{domandabox}

Una formula è una \textbf{tautologia} (o formula valida) se è vera in \emph{tutti} i modelli possibili, indipendentemente da come si assegnano i valori di verità alle variabili. L'esempio classico è $A \lor \neg A$: qualunque valore abbia $A$, la disgiunzione è sempre vera. Le tautologie sono importanti perché il \textbf{teorema di deduzione} afferma che $KB \models \alpha$ se e solo se $(KB \Rightarrow \alpha)$ è una tautologia.

Una \textbf{contraddizione} (formula insoddisfacibile) è l'opposto: è falsa in tutti i modelli. $A \land \neg A$ è una contraddizione. Questo concetto è alla base della dimostrazione per refutazione: se $KB \land \neg\alpha$ è una contraddizione, allora $KB \models \alpha$.

% =====================================================================
\section{Machine Learning}
% =====================================================================

\begin{domandabox}{Come si costruisce un albero di decisione? Cos'è l'algoritmo di Hunt?}
\end{domandabox}

L'\textbf{algoritmo di Hunt} è il nome del metodo ricorsivo classico per costruire alberi di decisione, da cui derivano algoritmi più noti come ID3 e C4.5. La costruzione di un albero di decisione segue un approccio ricorsivo. Si parte dall'intero insieme di training e si verifica la condizione di terminazione: se tutte le istanze appartengono alla stessa classe, si crea una foglia con quella classe. Se il nodo contiene poche istanze (sotto una soglia di pre-pruning) si crea una foglia con la classe maggioritaria. Altrimenti, si sceglie l'attributo che produce lo split migliore.

La bontà di uno split si misura con il \textbf{guadagno di informazione}: si calcola l'entropia del nodo padre, poi la media pesata delle entropie dei nodi figli, e si sceglie l'attributo che massimizza la differenza. Si creano quindi tanti nodi figli quanti sono i valori dell'attributo scelto, si partizionano le istanze tra di essi, e si applica la procedura ricorsivamente.

È essenziale distinguere il criterio di terminazione: il \textbf{numero di istanze sotto una soglia} è un valido criterio di pre-pruning. Non lo è invece l'overfitting del nodo (non rilevabile durante la costruzione) né un valore della funzione di valutazione (quella appartiene a Minimax, non agli alberi di decisione).

\begin{esempiobox}
\textbf{Dataset Tennis}: attributi = Outlook, Temperatura, Umidità, Vento; classe = Gioco/Non gioco. \\
Al nodo radice: entropia = $H = -\frac{9}{14}\log_2\frac{9}{14} - \frac{5}{14}\log_2\frac{5}{14} \approx 0.94$. \\
Split su Outlook: ramo ``Sunny'' ha 2 Sì e 3 No (entropia 0.97), ramo ``Overcast'' ha 4 Sì (entropia 0), ramo ``Rain'' ha 3 Sì e 2 No (entropia 0.97). Gain(Outlook) $\approx 0.246$ -- il più alto, quindi Outlook è la radice.
\end{esempiobox}

\begin{ricordabox}
\textbf{Costruzione = scegli sempre l'attributo che ``pulisce'' di più i nodi} (massimizza l'information gain). \textbf{Terminazione = foglia quando nodo puro OPPURE poche istanze} (pre-pruning). Mai fermarsi ``perché c'è overfitting'': non lo puoi sapere durante la costruzione.
\end{ricordabox}

\begin{domandabox}{Cosa misura l'entropia? E l'Information Gain?}
\end{domandabox}

L'\textbf{entropia} misura il grado di \emph{impurezza} o disordine in un insieme di dati: ci dice quanto sono mescolate le classi in un nodo. La formula è:
\[
H(t) = -\sum_{i} p(i|t) \log_2 p(i|t)
\]
dove $p(i|t)$ è la proporzione di istanze della classe $i$ nel nodo $t$. Quando tutte le istanze appartengono alla stessa classe, l'entropia è zero (nodo puro). Quando le classi sono distribuite in modo uniforme, l'entropia è massima.

L'\textbf{information gain} misura invece la \emph{riduzione di entropia} prodotta da uno split sull'attributo $A$:
\[
\text{Gain}(A) = H(\text{padre}) - \sum_{v} \frac{|D_v|}{|D|} H(D_v)
\]
Si sceglie l'attributo con gain massimo perché è quello che riduce di più l'incertezza. Una cosa importante da ricordare è che né l'entropia né il guadagno di informazione si calcolano dalla matrice di confusione: appartengono alla fase di costruzione del modello sul training set, non alla fase di valutazione.

\begin{ricordabox}
\textbf{Entropia alta = disordine = non sai niente.} Lanciare una moneta: $H = 1$ bit (massima incertezza). Moneta che dà sempre testa: $H = 0$ (certezza assoluta).\\[3pt]
\textbf{Information gain = quanto ordine guadagni splittando.} Se dopo lo split tutti i nodi figli sono puri ($H=0$), il gain è uguale all'entropia del padre: hai guadagnato tutto. Se i figli sono mescolati quanto il padre, gain = 0: non hai imparato niente.
\end{ricordabox}

\begin{domandabox}{Cosa si intende per overfitting? Cosa sono pre-pruning e post-pruning?}
\end{domandabox}

L'\textbf{overfitting} si verifica quando un modello si adatta troppo perfettamente al training set, imparando anche il rumore e le irregolarità casuali dei dati, invece di cogliere il pattern sottostante. Il risultato è un'alta accuratezza sul training ma basse prestazioni sui dati nuovi. Il principio di Occam suggerisce che, a parità di errore sul training, si preferisca il modello più semplice.

Il \textbf{pre-pruning} è una strategia per prevenire l'overfitting durante la costruzione dell'albero: si interrompe la ramificazione prima del completamento, quando il numero di istanze nel nodo scende sotto una soglia o il guadagno di informazione è troppo basso. Il rischio è l'underfitting: fermarsi troppo presto.

Il \textbf{post-pruning} è invece applicato dopo aver costruito l'albero completo: si esaminano i sottoalberi e si sostituiscono quelli che generalizzano male -- verificato su un validation set -- con semplici foglie etichettate con la classe maggioritaria. È in generale più affidabile del pre-pruning perché si dispone di una visione globale dell'albero. Il contesto del pruning è \emph{esclusivamente} l'apprendimento automatico, non i CSP né i giochi.

\begin{ricordabox}
\textbf{Overfitting = ``sa la lezione a memoria ma non capisce''.} Un albero che impara ogni singolo esempio di training ha errore 0 sul training ma generalizza male -- come studiare solo gli esercizi svolti e sbagliare all'esame con varianti.\\[3pt]
\textbf{Pre-pruning = smetti di studiare in anticipo} (rischi di non aver capito abbastanza). \textbf{Post-pruning = studia tutto, poi dimentica i dettagli inutili} (più sicuro ma più costoso).
\end{ricordabox}

\begin{domandabox}{Cos'è il Perceptron? Quale problema lo ha reso famoso?}
\end{domandabox}

Il percettrone, proposto da Rosenblatt nel 1957, è il modello più semplice di neurone artificiale: riceve un vettore di input $\mathbf{x}$, calcola la combinazione lineare pesata $\mathbf{w} \cdot \mathbf{x} + b$ e applica una funzione di attivazione per produrre l'output. L'apprendimento avviene attraverso la \textbf{regola delta}: i pesi vengono aggiornati proporzionalmente all'errore commesso sull'istanza corrente:
\[
w_j(t+1) = w_j(t) + \eta \cdot (d - o) \cdot x_j
\]
Il percettrone implementa un \emph{test lineare}: i pesi definiscono un iperpiano nello spazio degli input, e le istanze vengono classificate in base a quale lato dell'iperpiano si trovano.

Il problema che ha messo in luce i suoi limiti è il \textbf{problema dell'OR esclusivo (XOR)}. Minsky e Papert dimostrarono nel 1969 che il percettrone non può imparare la funzione XOR perché le sue quattro combinazioni di input non sono linearmente separabili: non esiste nessun iperpiano che separi correttamente i casi positivi da quelli negativi. Questo portò al primo ``inverno dell'IA'', fino allo sviluppo del backpropagation per reti multistrato negli anni '80.

\begin{esempiobox}
XOR: (0,0)$\to$0, (0,1)$\to$1, (1,0)$\to$1, (1,1)$\to$0. Nel piano 2D i punti positivi sono (0,1) e (1,0); i negativi sono (0,0) e (1,1). Non esiste nessuna retta che li separi: i positivi sono agli angoli opposti di un quadrato. Il percettrone traccia solo rette, quindi fallisce.\\
AND invece: (0,0)$\to$0, (0,1)$\to$0, (1,0)$\to$0, (1,1)$\to$1. Bastano rette: la retta $x_1 + x_2 > 1.5$ separa perfettamente.
\end{esempiobox}

\begin{ricordabox}
\textbf{Percettrone = righello.} Se puoi separare i dati con un righello (una linea retta / iperpiano), il percettrone impara. Se i dati si intrecciano come XOR, nessun righello funziona. Soluzione: più strati (reti neurali multilayer) che creano confini curvi.
\end{ricordabox}

\begin{domandabox}{Cos'è K-NN? Che differenza c'è con gli algoritmi che inducono regole?}
\end{domandabox}

Il K-Nearest Neighbours è un classificatore \textbf{lazy} (pigro): non costruisce nessun modello esplicito durante la fase di training, ma si limita a memorizzare tutto il dataset. Quando arriva una nuova istanza da classificare, calcola la distanza da tutti gli esempi noti, identifica i $k$ più vicini e assegna all'istanza la classe maggioritaria tra questi $k$ vicini.

Questa è la differenza fondamentale rispetto agli algoritmi che \textbf{inducono regole}: algoritmi come ID3 (alberi di decisione), o le strategie Specific-to-General e General-to-Specific, costruiscono un modello generalizzato durante il training -- un albero, un insieme di regole, un iperpiano -- che poi applicano alle nuove istanze. K-NN invece non generalizza nulla: il ``modello'' è l'intero training set.

Le strategie \textbf{Specific-to-General} partono da ipotesi molto specifiche (es. una singola istanza positiva) e le generalizzano gradualmente fino a coprire tutti gli esempi positivi. Le strategie \textbf{General-to-Specific} fanno l'opposto: partono da ipotesi molto generali (che coprono tutto) e le specializzano per escludere gli esempi negativi. Entrambe inducono regole esplicite, cosa che K-NN non fa mai.

% =====================================================================
\section{Matrice di Confusione}
% =====================================================================

\begin{domandabox}{Cos'è la matrice di confusione? Cosa si calcola da essa?}
\end{domandabox}

La matrice di confusione è uno strumento di valutazione dei modelli di classificazione: mette a confronto le etichette reali con quelle predette dal modello. Per una classificazione binaria, le quattro celle contengono i \textbf{veri positivi} (TP: predetto positivo, realmente positivo), i \textbf{falsi positivi} (FP: predetto positivo, realmente negativo), i \textbf{falsi negativi} (FN) e i \textbf{veri negativi} (TN).

Una prestazione ottimale corrisponde a tutti i valori concentrati sulla \textbf{diagonale principale} -- ovvero TP e TN -- con zero fuori dalla diagonale: significa che ogni istanza è stata classificata correttamente.

Le metriche che si calcolano dalla matrice di confusione includono l'\textbf{accuratezza} $= (TP + TN) / \text{totale}$, l'error rate $= 1 - \text{accuratezza}$, la precisione e il richiamo. Ciò che \emph{non} si calcola dalla matrice è l'entropia, l'information gain o l'overfitting: queste grandezze appartengono alla costruzione del modello, non alla sua valutazione sulle predizioni.

\begin{esempiobox}
100 email: 90 spam, 10 normali. Un classificatore dice ``tutto spam'': TP=90, FP=10, FN=0, TN=0. Accuratezza = 90\% -- sembra ottimo! Ma precision = 90/100 = 90\%, recall = 90/90 = 100\%, ma TN = 0: non riconosce mai un'email normale. Diagonale principale: solo TP e TN, FP e FN tutti zero = classificatore perfetto.
\end{esempiobox}

\begin{ricordabox}
\textbf{Matrice di confusione = tabella ``previsto vs reale''.} Diagonale = indovinato, fuori diagonale = sbagliato. \textbf{FP = allarme falso} (dici positivo ma è negativo). \textbf{FN = mancato allarme} (dici negativo ma è positivo). In medicina preferisci FP a FN: meglio un esame in più che una malattia mancata.
\end{ricordabox}

% =====================================================================
\section{Ontologie}
% =====================================================================

\begin{domandabox}{Cosa si intende per ontologia? Quali relazioni contiene?}
\end{domandabox}

Un'ontologia è una specifica formale di una concettualizzazione: definisce quali oggetti esistono in un dominio e come sono collegati tra loro. Le relazioni fondamentali sono quattro.

La relazione \textbf{isa} (is-a) esprime una \textbf{sottoclasse}: ``virologo isa microbiologo'' significa che virologo è una sottoclasse di microbiologo, e quindi ogni virologo è anche un microbiologo. La relazione \textbf{part-of} esprime la \textbf{composizione}: ``capsomero part-of virus'' indica che il capsomero è una componente strutturale del virus. La relazione \textbf{member} esprime l'\textbf{istanza}: ``HIV member virus'' indica che HIV è un esemplare specifico della classe virus, non una sottoclasse. Infine, un'espressione del tipo ``member(X, C) $\Rightarrow$ P(X)'' definisce una \textbf{proprietà di classe}: ogni membro di C ha la proprietà P.

Un insieme di sottocategorie che sono sia \emph{disgiunte} che \emph{esaustive} rispetto alla categoria padre costituisce una \textbf{partizione}. In una tassonomia, la proprietà corretta è che \emph{le sottocategorie di ciascuna categoria} formino una partizione, non che l'intero insieme di tutte le categorie lo sia.

RDF, il Resource Description Framework alla base del Web Semantico, rappresenta la conoscenza come \textbf{triple $\langle$soggetto, predicato, oggetto$\rangle$} e non tramite clausole di Horn.

\begin{esempiobox}
\textbf{Ontologia medica}: Socrate \textbf{member} Filosofo (istanza). Filosofo \textbf{isa} Umano (sottoclasse). Cervello \textbf{part-of} Umano (composizione). member(X, Umano) $\Rightarrow$ Mortale(X) (proprietà di classe). $\{$Greco, Romano, Egiziano$\}$ rispetto ad Antico = \textbf{partizione} (disgiunti ed esaustivi).\\[3pt]
Conseguenza per ereditarietà: Socrate member Filosofo, Filosofo isa Umano $\Rightarrow$ Socrate è Umano $\Rightarrow$ Socrate è Mortale.
\end{esempiobox}

\begin{ricordabox}
\textbf{isa = ``è un tipo di''} (Filosofo è un tipo di Umano -- vale per tutta la categoria). \textbf{member = ``è un individuo di''} (Socrate è un individuo della categoria Filosofo). \textbf{part-of = ``è un pezzo di''} (il cervello è un pezzo dell'umano). Confondere isa con member è l'errore tipico: Socrate non è una \emph{categoria} di Filosofo, è un \emph{membro}.
\end{ricordabox}

% =====================================================================
\section{Frame Problem e Situation Calculus}
% =====================================================================

\begin{domandabox}{Cos'è il Frame Problem?}
\end{domandabox}

Il frame problem è una delle sfide più sottili nella rappresentazione della conoscenza per agenti che agiscono nel mondo: è l'incapacità di sapere quali proprietà di una situazione si \textbf{preservano} dopo l'applicazione di un'azione. Il problema non è descrivere cosa cambia -- per quello bastano gli assiomi degli effetti -- ma descrivere tutto ciò che \emph{non} cambia senza dover elencare esplicitamente ogni singola proprietà invariata.

Per esempio, se un robot sposta un blocco, intuitivamente tutto il resto del mondo rimane uguale. Ma in una rappresentazione logica formale, come faccio a sapere che il colore del muro non è cambiato? Senza un meccanismo apposito, dovrei scrivere un assioma per ogni proprietà del mondo che non è stata modificata.

Nel \textbf{situation calculus} questo si risolve con gli \textbf{assiomi del successore di stato}: un fluente è vero nella situazione risultante dall'azione se e solo se l'azione lo ha reso vero, oppure era già vero e l'azione non lo ha reso falso. Questi assiomi catturano in modo compatto sia gli effetti sia la persistenza.

La definizione di un'azione nel situation calculus include le \textbf{precondizioni} (le proprietà che devono valere per applicare l'azione) e gli \textbf{effetti} (le proprietà che l'azione modifica). Le ``proprietà atemporali'' non fanno parte di questa definizione: sono assiomi della KB generale, non specifici di una singola azione.

\begin{esempiobox}
Azione \textsc{Sposta}(blocco, da, a): precondizione = blocco è in posizione ``da'' e ``a'' è libera; effetto positivo = blocco è in ``a''; effetto negativo = blocco non è più in ``da''. Assioma del successore di stato: On(b, a, Do(\textsc{Sposta}(b,x,a), S)) $\Leftrightarrow$ [l'azione l'ha reso vero] $\lor$ [era già vero e l'azione non l'ha reso falso].
\end{esempiobox}

\begin{ricordabox}
\textbf{Frame problem = il problema del ``tutto il resto rimane uguale''.} Quando sposti un blocco, il colore del soffitto non cambia -- ma devi dirlo esplicitamente in logica! Gli assiomi del successore di stato del situation calculus risolvono questo: ``vero dopo = l'azione l'ha reso vero OPPURE era già vero e l'azione non l'ha cancellato''. Questo cattura sia gli effetti sia la persistenza in una sola formula.
\end{ricordabox}

\end{document}
